\section{Conclusiones}

Con 374 módulos se obtiene una potencia instalada de \SI{168.3}{kWp} y una generación estimada equivalente al 98.4\,\% del consumo anual del edificio. El arreglo ocupa cerca del 54\,\% del techo considerado, así que hay espacio para la distribución propuesta. Las 22 cadenas de 17 módulos se mantienen dentro de los límites eléctricos de los inversores en las temperaturas evaluadas.

Para la acometida de 120/\SI{208}{V} se seleccionaron cinco inversores de \SI{25}{kW} fabricados para esa tensión. El quinto inversor encarece el sistema: la alternativa de \SI{153}{kWp}, con cuatro inversores, tiene un VPN de ¢34.4 millones frente a ¢29.5 millones del sistema seleccionado, aunque su generación cubre solo el 89.3\,\% del consumo.

Con el perfil de consumo supuesto, la energía comprada a la red disminuiría cerca de un 97\,\%. La factura, en cambio, se reduciría un 66.4\,\%, principalmente porque el cargo por demanda se mantiene y se agrega la tarifa TDER. Para saber si el sistema también reduciría la demanda facturada se necesitan mediciones horarias del consumo.

Con el precio estimado y el pago de contado, la inversión se recuperaría en 9.2 años, con una TIR de 9.7\,\%, mayor que la tasa de descuento de 7\,\%. Si se financia el 80\,\% a diez años, la cuota sería mayor que el ahorro neto durante el plazo del préstamo, en unos ¢1.3 millones por año. La cooperativa tendría que cubrir esa diferencia hasta terminar de pagarlo.

La rentabilidad depende de que se mantengan las condiciones actuales de la medición neta sencilla, porque cerca de la mitad de la generación estimada se deposita en la red. Si esa energía se vendiera como excedente o no se reconociera, el VPN sería negativo. Antes de tomar una decisión de inversión deben confirmarse las condiciones de medición neta con el ICE, el precio mediante cotizaciones, la capacidad estructural del techo y las condiciones de conexión tanto en el tablero principal como en la red de distribución.
